\section{Current claims and their dependencies}\label{sec:claims}
``Written proof'' below describes the arguments in this manuscript.
It does not mean external peer review or formal verification. The
coefficient-prime restrictions are part of the conclusions.

\small
\begin{longtable}{>{\raggedright\arraybackslash}p{0.23\textwidth}>{\raggedright\arraybackslash}p{0.35\textwidth}>{\raggedright\arraybackslash}p{0.32\textwidth}}
\toprule
Conclusion & Proof location and main inputs & Scope or remaining limit\\
\midrule
\endfirsthead
\toprule Conclusion & Proof location and main inputs & Scope or remaining limit\\
\midrule\endhead
Ranks and rational descents & Sections~\ref{sec:geometry}, \ref{sec:descent}:
parabolic dimension, literal inversion and deck projectors &
Specified arithmetic normalization; $W_4^-=W_4^+\otimes\chi_{-3}$.\\[4pt]
Pairing and good Euler factors & Sections~\ref{sec:descent}, \ref{sec:inert}:
cup product, exact finite sums and inert boundary term &
Exact formulas at $p\ge5$; no certification of every stored numerical row.\\[4pt]
Good-prime weight five & Corollary~\ref{cor:descent-good-prime-purity}:
Deligne's weight theorem and specialization &
Independent of empirical reciprocal-quartic tests.\\[4pt]
Rational Chow realization and Hodge numbers &
Theorem~\ref{thm:scholl-realization}: full-level cover, finite projectors,
Scholl comparison & All $\ell$; $h^{5,0}=h^{0,5}=2$; no other Hodge types.\\[4pt]
Local data at $3$ & Section~\ref{sec:local-three}: quotient family,
collision cohomology, ordered Frobenius flag &
Every $\ell\ne3$; plus $P_3=1-3T$, $a_3=3$; minus $P_3=1$, $a_3=4$.\\[4pt]
Local data at $2$ & Section~\ref{sec:local-two}: exact Galois closure,
semistable charts, tame comparison, tails and nodes &
Every $\ell\ne2$; both $d_2=0$, $P_2=1$, $\Sw_2=0$, $a_2=4$.\\[4pt]
Conductor products & Section~\ref{sec:conductors}: both local
calculations and good-prime unramifiedness &
Plus $432$, minus $1296$; for $\ell\ne2,3$, these are the prime-to-$\ell$
Artin conductors.\\[4pt]
Strong absolute irreducibility & Section~\ref{sec:irreducibility}:
single-block monodromy at $3$, Hodge--Tate exponents and rank-one
characters over $\Q$ & Every $\ell\ne3$, after every finite
number-field extension; coefficient prime $3$ is not covered.\\[4pt]
Cubic second-moment comparison & Section~\ref{sec:cubic}:
geometric trace/determinant, stated comparison theorem, seven exact rows
and Sturm bound $24$ & Separate rank-two object; bounded exact
certificate retained; orders $4,6$ do not inherit this proof.\\
\bottomrule
\end{longtable}
\normalsize

The local factor and exponent independence asserted here is obtained
directly for all allowable auxiliary primes. A full comparison of the
bad-prime Weil--Deligne representations, including a treatment at equal
coefficient and residue characteristic, is a different statement.
The remaining research questions are listed in Section~\ref{sec:limits}.

\paragraph{Sign dictionary.}
At an inert good prime, the literal plus descent $I=+1$ has trace
$-S_p-1$ and conductor $432$. The literal minus descent $I=-1$ has trace
$S_p+1$ and conductor $1296$. Thus an older numerical label ``plus sign''
for the finite-sum formula means the latter object; it is not the
literal plus eigenspace of $I$.

\paragraph{Review priorities.}
The most useful independent review targets are the arithmetic constant
in the dyadic closure, the coefficient action on the main component and
nodes, the horizontal-boundary nearby-cycle hypotheses, and the
local-at-$3$ global-to-local map with its Frobenius normalization.
The strong-irreducibility deduction depends on the complete local flag,
not merely on its conductor exponent. These are review targets within
the stated proof, not newly claimed external endorsements.
